%% kapa: exam=1 round=25 year=2014 period=2 form=A subject=appraisal_law subject_ko=부동산관계법규 seq=33 number=33
%% source: Q-Net 감정평가사 기출문제 (cst00309) artlSeq=5101207 file=extracted/1008/5101207/2014년 제25회 감정평가사 1차시험 기출문제/2교시(A형).hwp
%% conversion: kordoc 4.24.51 -> tools/gen1.py (auto); review=proofread-sonnet; flags=none
\begin{question}{33}
측량ㆍ수로조사 및 지적에 관한 법령상 토지의 용도와 지목이 옳게 연결된 것은?
\begin{choices}
  \item 물을 상시적으로 이용하지 않고 과수류를 집단적으로 재배하는 토지 - 전
  \item 묘지의 관리를 위한 건축물의 부지 - 대
  \item 식용으로 죽순을 재배하는 토지 - 과수원
  \item 자동차정비공장 내 급유시설부지 - 주유소용지
  \item 종교용지로 된 토지에 있는 유적을 보호하기 위하여 구획된 토지 - 사적지
\end{choices}
\end{question}
\answer{2}
